%%%%%%%%%%%%%%%%%%%%%%%%%%%%%%%%%%%%%%%%%%%%%%%%%%%%%%%%%%%%%%%%%%%%%%%%%%%%%%%%
% Recent Grad Resume Template
% CV content migrated from sample.tex (AltaCV)
%%%%%%%%%%%%%%%%%%%%%%%%%%%%%%%%%%%%%%%%%%%%%%%%%%%%%%%%%%%%%%%%%%%%%%%%%%%%%%%%

\documentclass[margin]{res}
\newsectionwidth{1.45in}

\makeatletter
\@fileswtrue
\makeatother

\usepackage[utf8]{inputenc}
\usepackage[T5]{fontenc}
\usepackage[english]{babel}
\usepackage{helvet}
\usepackage{hyperref}
\usepackage{csquotes}
\DeclareQuoteGlyph{T5}{34}
\DeclareQuoteGlyph{T5}{39}
\usepackage[backend=biber,style=numeric]{biblatex}
\addbibresource{../public.bib}

\textheight=700pt

\makeatletter
\renewcommand{\namefont}{\LARGE\bfseries}
\newcommand{\lastupdate}{June 25, 2026}
\def\print@name{%
  \begingroup
  \leavevmode\hbox to \textwidth{%
    {\namefont\@name}%
    \hfil
    {\small Last update: \lastupdate}%
  }%
  \par\vskip 2pt
  \fullline
  \vskip 0.75em
  \endgroup
}
\makeatother

\begin{document}

\name{Nguyễn Chí Bằng}

\begin{resume}

  \section{CONTACT\\INFORMATION}
  \begin{tabular}[t]{@{}l@{\quad}l@{}}
    \textbf{Phone:} & \href{tel:+84832946009}{(+84) 832 946 009} \\
%   \textbf{University Email:} & \href{mailto:}{} \\
    \textbf{Personal Email:} & \href{mailto:chibangn1@gmail.com}{chibangn1@gmail.com} \\
    \textbf{Homepage:} & \href{https://bancie.github.io/}{https://bancie.github.io/}
  \end{tabular}



  \begin{format}
    \title{l}\employer{r}\\
    \dates{l}\location{r}\\
    \body\\
  \end{format}

  \section{EXPERIENCE}
  
  % \employer{7/2024 -- Present}
  % \location{}
  % \dates{}
  % \title{\href{https://www.vlu.edu.vn/institute/vien-khoa-hoc-tinh-toan-va-tri-tue-nhan-tao}{\textbf{Computational Science and Artificial Intelligence Research Institute (COSARI), VLU, Vietnam}}}

  \employer{Collaborative Researcher}
  \location{Van Lang University, Vietnam}
  \dates{7/2024 -- Present}
  \title{\href{https://www.vlu.edu.vn/institute/vien-khoa-hoc-tinh-toan-va-tri-tue-nhan-tao}{\textbf{Computational Science and Artificial Intelligence Research Institute (COSARI)}}}
  \begin{position}
    Built data-driven Python tools that combine time-series data and machine learning to optimize scheduling and logistics decisions.
    \begin{itemize}
      \item Applied time-series analysis and ML models to single/multi-machine scheduling and optimal facility location problems.
      \item Collaborative research with Ho Chi Minh City University of Technology (HCMUT).
      \item Supervisors: Dr. Lê Minh Huy (VLU), Dr. Trần Thanh Hiệp (HCMUT)
    \end{itemize}
  \end{position}

  \employer{Data Science Intern}
  \location{HCMUT, Vietnam}
  \dates{12/2025 -- 4/2026}
  \title{\href{http://imacs.hcmut.edu.vn}{\textbf{Institute of Mathematical and Computational Sciences (IMACS)}}}
  \begin{position}
    Built a vision-based UAV self-positioning pipeline for GPS-denied environments by matching drone images to satellite references.
    \begin{itemize}
      \item Developed and evaluated a DenseUAV baseline (ViT image feature extractor) for cross-view localization on UAV and satellite imagery.
      \item Trained models on cloud GPUs using Modal, with experiment tracking on Weights \& Biases (W\&B).
      \item Evaluation: Excellent.
      \item Supervisor: M.Sc. Nguyễn Văn Gia Thịnh
      \item \href{https://wandb.ai/chibangn1-https-bancie-github-io-tilearn-/denseuav-baseline/reports/DenseUAV-Non-GPS-Training-Performance--VmlldzoxNjE5NDMwOQ?accessToken=hl4axuavbu3up6rygu5oy7r3ru3jiwwlg285ax59kdhonfcqgx2n0nhmwex4rfob}{W\&B Report}: Training metrics and loss curves.
      \item \href{https://huggingface.co/Bancie/UAV-Self-Positioning-23M-ZCN}{Hugging Face}: Published model checkpoint.
      \item \href{https://github.com/Bancie/denseUAV_baseline}{GitHub}: Baseline model development code.
    \end{itemize}
  \end{position}

\section{COMPETITIONS}
\par
\href{https://www.kaggle.com/competitions/bank-churn-competition-by-m-linside-ai-specialization}{\textbf{Bank Churn Challenge by MLinside AI Specialization}} (6/2026): Kaggle competition on bank customer churn prediction.
\begin{itemize}
  \item Predicted the probability that a bank customer will close their account (binary classification; target: \texttt{Exited}).
  \item Built an ML classification pipeline with customer-feature engineering and model evaluation using ROC-AUC.
  \item Ranked \textbf{5th} on the public leaderboard (Score (ROC-AUC): 0.93028).
  \item \href{https://www.kaggle.com/competitions/bank-churn-competition-by-m-linside-ai-specialization}{Link Kaggle Competition}
  \item Kaggle Notebook: \textit{link forthcoming} % TODO: \href{https://www.kaggle.com/code/...}{Notebook}
\end{itemize}

\par
\href{https://sgu.edu.vn}{\textbf{Methods for Integer Linear Programming Problems}} (5/2024): University-level scientific research project at Saigon University, Vietnam.
\begin{itemize}
  \item Applied machine learning to process experimental data for evaluating Branch and Bound and Gomory methods for integer linear programming.
  \item Grade: Excellent.
  \item Supervisor: Assoc. Prof. Dr. Tạ Quang Sơn
\end{itemize}

  \section{PROJECTS}
  
  \par
  \href{https://bancie.github.io/TiLearn/}{\textbf{TiLearn}} (7/2024 -- Present): An open-source data science platform for task priority and duration prediction.
  \begin{itemize}
    \item Used time-series, work, and personal data to train ML models that predict and rank task priority and estimated execution time.
    \item Built a Python library and open-source platform for integrating predictions into time-management workflows.
    \item Supervisors: Dr. Lê Minh Huy (VLU), Dr. Trần Thanh Hiệp (HCMUT)
    \item \href{https://pypi.org/project/TiLearn/}{PyPI}
    \item \href{https://github.com/Bancie/TiLearn}{GitHub}
  \end{itemize}

  
  \par
  \href{https://github.com/Bancie/LearningDB}{\textbf{LearningDB}} (6/2024 -- Present): A personal data platform for habit tracking, ML pipelines, and natural-language daily planning.
  \begin{itemize}
    \item Built a habit-tracking system to collect structured personal data and run ML models for data mining.
    \item Exposed habit and personal data via MCP tools and integrated a LangChain-based conversational assistant for day-to-day planning and task management.
    \item Deployed as a multi-service stack (API, orchestrator, web) with Docker Hub images.
    \item \href{https://github.com/Bancie/LearningDB}{GitHub}
  \end{itemize}

  \par
  \href{https://github.com/Bancie/scheloc}{\textbf{ScheLoc}} (11/2025 -- Present): Optimize transport time using geographic and time-series data.
  \begin{itemize}
    \item Combined geographic data with time-series data to optimize transport and delivery schedules.
    \item Discrete Optimization Research Group, Can Tho University, Vietnam
    \item \href{https://github.com/Bancie/scheloc}{GitHub}
  \end{itemize}


  \section{VOLUNTEERS}
  \par
  \href{https://github.com/blazickjp/arxiv-mcp-server}{\textbf{ArXiv MCP Server}} (4/2026 -- 7/2026): Open-source contribution to an MCP server that enables AI assistants to search and analyze arXiv papers through a simple Model Context Protocol interface.
  \begin{itemize}
    \item Fixed truncated arXiv PDF downloads by streaming \texttt{pdf\_url} with httpx (merged \href{https://github.com/blazickjp/arxiv-mcp-server/pull/99}{PR \#99}; related \href{https://github.com/blazickjp/arxiv-mcp-server/issues/98}{issue \#98}).
    \item Ranked top 2 human contributor (excluding bots) over Apr--Jul 2026.
    \item \href{https://github.com/blazickjp/arxiv-mcp-server}{GitHub}
    \item \href{https://mcpmarket.com/server/arxiv-1}{MCP Market}
  \end{itemize}

  \par
  \textbf{Reading Performance} (7/2025 -- Present): A personal system to collect essential life data.
  \begin{itemize}
    \item Developed a system for collecting and analyzing personal data.
    \item Reading performance dataset with demographic and contextual factors reached 40+ upvotes, 1,500 downloads, 5,500 views, and numerous analyses.
    \item \href{https://www.kaggle.com/datasets/chibangng/reading-performance}{Kaggle}
    \item \href{https://huggingface.co/datasets/Bancie/Reading-Dataset}{Hugging Face}
  \end{itemize}

\section{RESEARCH}
\nocite{*}
\printbibliography[heading=none]

  \section{EDUCATION}
  \textbf{Saigon University}, Ho Chi Minh City, Vietnam \\
  {\sl Bachelor's in Applied Mathematics and Computer Science}, 2022 -- 2026\\
  Grade: Good

  \section{ENGLISH}
  \textbf{VSTEP} - Overall: B2 (CEFR B2), 2026

  \section{STRENGTHS}
  \textbf{Soft skills}: Hard-working, Fast adaptability, Rapid reading, deep comprehension, Motivator \& Leader.
  \\
  \textbf{AI/ML}: AI, ML/DL, Data Analysis.
  \\
  \textbf{Programming}: C/C++, Python, FastAPI, PyTorch, TensorFlow, ...
  \\
  \textbf{Domain}: Optimization, Statistical, Operations Research.

  \section{REFERENCES}
  \textbf{Assoc. Prof. Dr. Tạ Quang Sơn}\\
  Faculty of Mathematics and Applications\\
  \href{mailto:taquangson@sgu.edu.vn}{taquangson@sgu.edu.vn}\\
  Saigon University, Ho Chi Minh City, Vietnam

  \par
  \textbf{Dr. Lê Minh Huy}\\
  Faculty of Basic Sciences\\
  \href{mailto:huy.lm@vlu.edu.vn}{huy.lm@vlu.edu.vn}\\
  Van Lang University, Ho Chi Minh City, Vietnam

  \par
  \textbf{Dr. Trần Thanh Hiệp}\\
  Faculty of Applied Sciences\\
  \href{mailto:ttthiep.sdh241@hcmut.edu.vn}{ttthiep.sdh241@hcmut.edu.vn}\\
  Ho Chi Minh City University of Technology, Vietnam
  
  \par
  \textbf{M.Sc. Nguyễn Văn Gia Thịnh}\\
  Institute of Mathematical and Computational Sciences (IMACS)\\
  \href{mailto:nguyenvangiathinh@hcmut.edu.vn}{nguyenvangiathinh@hcmut.edu.vn}\\
  Ho Chi Minh City University of Technology, Vietnam


\end{resume}
\end{document}
